% ============================================================
% TRACK 2: PIPELINE-GENERATED DOMAIN VALIDATION
% New addition — April 2026. For evaluation by Dan McCreary.
% ============================================================

\section{Track 2: Pipeline-Generated Domain Validation}\label{sec:track2}

\subsection{Motivation}

All 44 domains in the primary benchmark (Track 1) originate from the
McCreary Intelligent Textbook Corpus: hand-curated educational DAGs
where a domain expert manually mapped concept dependencies.
A critical open question is whether the CKG architecture's performance
advantage depends on that curation quality or whether it transfers to
\textit{programmatically constructed} knowledge graphs derived from
external data sources.

Track 2 answers this question by constructing a complete CKG domain
from scratch using no pre-existing expert-curated graph, no educational
corpus, and no McCreary source material. The domain selected is
\textbf{GLP-1/Obesity pharmacology}---a commercially active life
sciences domain with rapidly evolving clinical trial data, 8 FDA-approved
agents, and a pipeline of 150+ ongoing trials at the time of corpus
construction (April 2026).

\subsection{Data Source and Pipeline}

Track 2 uses ClinicalTrials.gov as the sole external data source,
accessed via the NIH/NLM API (v2). The pipeline consists of four stages:

\begin{enumerate}[leftmargin=2em]
  \item \textbf{API extraction.} Structured query against ClinicalTrials.gov
        returns 668 semaglutide trials, 224 tirzepatide trials, and 158
        pipeline agent trials (retatrutide, cagrisema, orforglipron, mazdutide).
        Trial metadata includes NCT identifiers, phase, enrollment, endpoints,
        mechanisms, and completion dates.

  \item \textbf{Concept extraction.} Trial data is parsed to extract
        pharmacological entities (agents, mechanisms, indications, trial
        programs, outcomes) and their dependency relationships (e.g.,
        \textit{SELECT trial} $\to$ \textit{semaglutide} $\to$
        \textit{GLP-1 receptor agonism}). Taxonomy labels are assigned
        from a domain-specific schema: \texttt{FOUND} (foundational
        mechanism), \texttt{DRUG} (approved agent), \texttt{TRIAL}
        (completed landmark trial), \texttt{PATH} (pathway/mechanism),
        \texttt{COMPL} (complication/adverse effect), \texttt{SPEC}
        (special population), \texttt{COMBO} (combination strategy).

  \item \textbf{Graph construction.} Extracted concepts and dependencies
        are written to \texttt{learning-graph.csv} in the standard CKG
        schema (\texttt{ConceptID, ConceptLabel, Dependencies, TaxonomyID}).
        The resulting graph contains \textbf{90 concepts} and \textbf{170
        dependency edges} covering foundational mechanisms through
        next-generation pipeline agents and cross-indication indications
        (cardiovascular, renal, neurological, addiction).

  \item \textbf{Query generation.} The standard benchmark harness
        (\texttt{generate\_queries.py}) runs unchanged on the
        \texttt{learning-graph.csv}, producing 170 queries in the T1--T5
        taxonomy with deterministic ground truth derived from DAG edges.
\end{enumerate}

The full pipeline---from raw API data to benchmark-ready domain---requires
no manual annotation and no subject matter expert review beyond the
initial taxonomy schema. We refer to this four-stage procedure as an
\textit{automated CKG construction pipeline} and report it here as a
research contribution: a reproducible method for generating
benchmark-ready Compact Knowledge Graph domains directly from
structured external data sources.

\subsection{Corpus Construction for RAG and GraphRAG}

To enable a fair three-way comparison, a prose corpus was constructed
from the same ClinicalTrials.gov data used to build the CKG. Five
structured narrative documents were written covering: (1) market overview
and approved agents, (2) landmark clinical trial evidence (STEP,
SURMOUNT, SELECT, SUMMIT, FLOW programs), (3) next-generation pipeline
intelligence, (4) indication expansion across 15+ disease areas, and
(5) investment-relevant signals including oral formulation, muscle
preservation, and CNS/addiction frontiers. These documents were indexed
by RAG (FAISS, 512-token chunks, top-5 retrieval) and GraphRAG (Microsoft
GraphRAG v1.x, local community reports) using the same configurations
as Track 1.

\subsection{Results}

Table~\ref{tab:track2-results} presents macro-average token-level F1
for all three systems on the GLP-1/Obesity domain alongside the
Track 1 benchmark aggregate for reference.

\begin{table}[htbp]
\centering
\caption{Track 2 results: pipeline-generated GLP-1/Obesity domain (170 queries)
compared with Track 1 benchmark aggregate (44 hand-curated domains, 7,758 queries).}
\label{tab:track2-results}
\begin{tabular}{llccccc}
\toprule
\textbf{Track} & \textbf{System} & \textbf{Macro F1} & \textbf{Tokens/q} &
  \textbf{Ret.\ Tokens/q} & \textbf{RDS} & \textbf{n queries} \\
\midrule
Track 1 & CKG      & 0.4709 & 269  & 44   & 0.00201 & 7,758 \\
Track 1 & RAG      & 0.1231 & 2,982 & 2,392 & 0.0000482 & 7,191 \\
Track 1 & GraphRAG & 0.1200 & 3,450 & ---  & 0.0000452 & 2,683 \\
\midrule
Track 2 & CKG      & \textbf{0.5298} & 346  & 54   & \textbf{0.00153} & 170 \\
Track 2 & RAG      & 0.1538 & 2,828 & 2,214 & 0.0000544 & 170 \\
Track 2 & GraphRAG & 0.1436 & 3,450 & ---  & 0.0000416 & 170 \\
\bottomrule
\end{tabular}
\end{table}

The Track 2 CKG F1 of \textbf{0.5298} exceeds the Track 1 CKG
macro-average of 0.4709 by 12.5\%. This result is notable because the
GLP-1 graph was generated by pipeline from raw API data, not curated by
a domain expert. The performance advantage is not degraded by automated
construction---it is \textit{maintained or improved}.

Table~\ref{tab:track2-by-type} shows F1 by query type for the
GLP-1/Obesity domain.

\begin{table}[htbp]
\centering
\caption{Track 2 F1 by query type --- GLP-1/Obesity domain (170 queries).}
\label{tab:track2-by-type}
\begin{tabular}{lccccc}
\toprule
\textbf{System} & \textbf{T1 entity} & \textbf{T2 dep.} & \textbf{T3 path} &
  \textbf{T4 aggr.} & \textbf{T5 cross} \\
\midrule
CKG      & \textbf{0.225} & \textbf{0.677} & \textbf{0.873} & \textbf{0.998} & \textbf{0.425} \\
RAG      & 0.150 & 0.076 & 0.221 & 0.108 & 0.226 \\
GraphRAG & 0.129 & 0.051 & 0.216 & 0.031 & 0.258 \\
\bottomrule
\end{tabular}
\end{table}

Several findings are noteworthy. \textbf{T4 (category aggregation)} reaches
CKG F1 $= 0.998$---near-perfect enumeration of agents by drug class,
indication by anatomy, and trial by program---compared with RAG's 0.108
and GraphRAG's 0.031. This is the sharpest divergence observed in either
track and confirms that taxonomy-indexed graphs enable exact-set retrieval
regardless of domain.

\textbf{T3 (multi-hop path)} reaches CKG F1 $= 0.873$, substantially above
the Track 1 T3 average of 0.660. The GLP-1 dependency graph encodes
mechanistic chains (receptor $\to$ signaling pathway $\to$ downstream
effect $\to$ clinical outcome) that BFS traversal resolves with high
precision. Automated extraction preserved these chains because the
source API data is itself structured.

\textbf{GraphRAG} scores slightly above RAG on T5 cross-concept queries
(0.258 vs.\ 0.226), consistent with its community-report architecture
capturing co-occurrence relationships. This mirrors Track 1 behavior.

\subsection{Implications for the CKG Architecture}

Track 2 establishes three findings beyond Track 1.

\paragraph{Finding 1: Retrieval performance depends on graph structure, not curation source.}
CKG F1 on the pipeline-generated GLP-1 domain (0.530) exceeds the
hand-curated educational average (0.471). If expert curation quality were
the driver of CKG performance, the pipeline domain would score lower. It
does not. This means the architecture generalizes: any domain with
stable concept relationships expressible in a DAG benefits from CKG
retrieval regardless of how the DAG was built.

\paragraph{Finding 2: The CKG factory is viable for commercial domains.}
The GLP-1 benchmark demonstrates an end-to-end automated pipeline
from public API data to a benchmarked knowledge retrieval system. The
pipeline requires no annotation budget, no expert review, and no existing
textbook or corpus---only a structured data source and a taxonomy schema.
This generalizes the CKG architecture beyond educational settings into
any knowledge-intensive commercial domain (pharmaceutical, legal, financial,
regulatory) where public or proprietary structured data is available.

\paragraph{Finding 3: The 28$\times$ RDS advantage holds on enterprise domains.}
RDS for CKG on GLP-1 is $0.00153$ versus RAG's $0.0000544$---a ratio of
28$\times$. The Track 1 RDS ratio is 42$\times$. The slight reduction on
Track 2 reflects the GLP-1 domain's higher prose complexity (longer
ground-truth answers for cross-indication queries) but the order-of-magnitude
token efficiency advantage is preserved. A life sciences organization
deploying CKG for structured pharmacology queries over RAG realizes
approximately 97\% cost reduction with $3.4\times$ accuracy improvement.

\subsection{Track 2 and the CKG Factory Paradigm}

Track 2 clarifies three implications of the benchmark results that bear
directly on how Compact Knowledge Graphs can be built and deployed in
practice.

\paragraph{Automated construction invalidates the ``expert curation required'' limitation.}
A recurring critique of structured-knowledge approaches is that they
depend on costly human curation and therefore do not scale. The Track 2
result is a direct empirical counterexample: a CKG produced by a
four-stage pipeline over a public API, with no expert-curated edges
and no textbook source material, matches and slightly exceeds the
retrieval performance of the 44 hand-curated Track 1 domains
(F1 $=$ 0.5298 vs.\ 0.4709). Expert curation is therefore not a
precondition for the CKG structural advantage; it is one of several
viable construction strategies.

\paragraph{Multi-domain ensemble demonstrates educational-to-commercial transfer.}
The Track 2 GLP-1/Obesity domain coexists with the 44 Track 1
educational domains in a single evaluation harness (\texttt{ckg\_harness.py}),
queried and scored identically. This heterogeneous ensemble---hand-curated
educational DAGs alongside a pipeline-generated pharmacology DAG---shows
that the CKG retrieval advantage is a property of the \emph{representation},
not of the domain. The same BFS/DFS subgraph-extraction procedure that
performs on calculus or microbiology performs on clinical trial
pharmacology, with no modification to the retrieval code and no
domain-specific tuning.

\paragraph{Implications for practitioners.}
For organizations considering CKG deployment in a commercial domain, the
Track 2 evidence suggests a concrete recipe: (i) identify a structured
data source whose records encode entities and dependencies (public APIs,
regulatory registries, internal knowledge bases, product catalogues);
(ii) define a lightweight taxonomy schema for the domain;
(iii) run the four-stage pipeline to produce a \texttt{learning-graph.csv};
(iv) query via the standard harness. This recipe is domain-agnostic,
requires no annotation budget, and produces a benchmarkable CKG with
retrieval performance comparable to hand-curated knowledge graphs in
the same architectural class.

The Track 2 data and corpus are available in the benchmark repository
under \texttt{benchmark/domains/glp1-obesity/} and
\texttt{corpus/glp1-obesity/}.
